\documentclass[12pt,a4paper]{article}
\usepackage[T2A]{fontenc}
\usepackage[utf8]{inputenc}
\usepackage[russian]{babel}
\usepackage{amsmath,amssymb}
\usepackage[margin=2cm]{geometry}
\usepackage{hyperref}
\usepackage{xcolor}

\definecolor{titleblue}{RGB}{26,35,126}
\definecolor{headblue}{RGB}{40,53,147}
\definecolor{subblue}{RGB}{57,73,171}

\title{\color{titleblue}\textbf{Математическое описание методов\\сглаживания сигналов}}
\author{}
\date{}

\begin{document}
\maketitle

%% === ВВЕДЕНИЕ ===
\section{Введение}

Рассмотрим дискретный сигнал $x[0], x[1], \ldots, x[N{-}1]$, представляющий собой последовательность измерений некоторой случайной величины. Предполагается, что каждое измерение состоит из истинного сигнала и аддитивной компоненты шума:
\[
x[i] = \mu[i] + \varepsilon[i],
\]
где $\mu[i]$ --- истинное значение в момент $i$, а $\varepsilon[i]$ --- случайная погрешность с нулевым математическим ожиданием и дисперсией $\sigma^2$.

Задача сглаживания состоит в восстановлении функции $\mu[i]$ по зашумлённым наблюдениям. С точки зрения математической статистики это задача \textbf{непараметрической регрессии}, в которой ищется оценка условного математического ожидания $\mathbb{E}[x \mid i] = \mu[i]$.

Ниже описаны четыре метода, каждый из которых может быть интерпретирован как особый вид \textbf{ядерного сглаживания} (kernel smoothing).

%% === 1. BOXCAR ===
\section{Равномерное скользящее среднее (Boxcar)}

\subsection{Статистическая модель}

Равномерное скользящее среднее является простейшей формой оценки условного математического ожидания. Оно основано на идее, что в окне фиксированного размера $M$ значения сигнала мало изменяются, и поэтому их среднее арифметическое является хорошей оценкой истинного уровня в центре окна.

С точки зрения \textbf{ядерной оценки плотности}, boxcar-фильтр использует равномерное (прямоугольное) ядро:
\[
K(u) = \frac{1}{2}, \quad |u| \leq 1; \qquad K(u) = 0, \quad |u| > 1.
\]
Полуширина окна $m$ играет роль параметра ширины (bandwidth) $h$, определяющего степень сглаживания.

\subsection{Свойства оценки}

Математическое ожидание оценки boxcar для точки $i$ равно среднему $\mu$ по окну вокруг $i$. Если истинная функция $\mu$ постоянна, оценка является \textbf{несмещённой}. Если $\mu$ изменяется внутри окна, возникает систематическая ошибка (bias), пропорциональная второй производной $\mu$ и квадрату полуширины окна:
\[
\operatorname{bias} = O(m^2 \cdot \mu'').
\]

Дисперсия оценки определяется дисперсией шума и размером окна:
\[
\operatorname{Var} = \frac{\sigma^2}{W(i)},
\]
где $W(i)$ --- число элементов в окне. На границах сигнала окно уменьшается, поэтому дисперсия растёт.

С точки зрения теории оценок, boxcar демонстрирует классический \textbf{компромисс между смещением и дисперсией} (bias--variance tradeoff). Увеличение окна уменьшает дисперсию, но увеличивает смещение, и наоборот. Оптимальный размер окна минимизирует среднеквадратичную ошибку:
\[
\operatorname{MSE} = \operatorname{bias}^2 + \operatorname{Var}.
\]

\subsection{Граничные условия}

На границах сигнала окно не может быть полностью заполнено. В данной реализации применяется \textbf{обрезка} (clamping): окно уменьшается, охватывая только имеющиеся элементы. Это приводит к увеличению дисперсии на границах, но сохраняет несмещённость, поскольку все используемые элементы являются подлинными наблюдениями.

%% === 2. GAUSSIAN ===
\section{Гауссов фильтр (Gaussian filter)}

\subsection{Ядерная оценка с гауссовским ядром}

Гауссов фильтр является классическим примером \textbf{ядерной оценки условного среднего} (Nadaraya--Watson estimator). Вместо равномерного веса каждому элементу окна присваивается вес, убывающий по закону Гаусса от центра окна к краям. Это соответствует использованию гауссовой ядерной функции:
\[
K(u) = \frac{1}{\sqrt{2\pi}} \exp\!\left(-\frac{u^2}{2}\right).
\]

Параметр $\sigma = M/6$ определяет ширину ядра и задаётся так, чтобы интервал $\pm 3\sigma$ укладывался в окно. Это означает, что более 99.7\% суммарного веса приходится на элементы внутри окна.

\subsection{Смещение и дисперсия}

Для гауссова ядра смещение оценки также пропорционально второй производной $\mu$ и квадрату параметра ширины: $\operatorname{bias} = O(\sigma^2 \cdot \mu'')$. Однако, в отличие от boxcar, гауссов фильтр обладает более мягким затуханием весов, что снижает эффект \textbf{артефактов Гиббса} --- нежелательных колебаний на частотах вблизи частоты среза фильтра.

Дисперсия гауссова фильтра аналогична boxcar: $\operatorname{Var} \sim \sigma^2 \cdot \int K^2 / (N \cdot h)$, где $\int K^2 = 1/(2\sqrt{\pi})$ --- интеграл квадрата ядра. Это несколько больше, чем для равномерного ядра ($\int K^2 = 1$), поэтому при прочих равных гауссов фильтр имеет незначительно бо́льшую дисперсию, но компенсирует это отсутствием артефактов.

\subsection{Симметричность и фазовые свойства}

Гауссово ядро является чётной функцией, что обеспечивает \textbf{нулевую фазовую характеристику} фильтра. Это означает, что фильтр не сдвигает сигнал во времени --- пики и переходы остаются на тех же позициях. Симметрия ядра гарантирует отсутствие фазовых искажений, что критически важно при обработке импульсных и знакопеременных сигналов.

\subsection{Граничные условия}

В отличие от boxcar, гауссов фильтр использует \textbf{симметрическое отражение} (mirror reflection) данных за пределами массива. С точки зрения теории оценок, это соответствует предположению о симметрии процесса за пределами наблюдаемой области. Отражение сохраняет полную ширину окна на границах, обеспечивая равномерную дисперсию по всему сигналу.

%% === 3. PROGRESSIVE ===
\section{Прогрессивное среднее (Progressive Mean)}

\subsection{Оценка с расширяющимся окном}

Прогрессивное среднее можно интерпретировать как последовательную (рекурсивную) оценку среднего уровня сигнала, в которой на каждом шаге используются все имеющиеся наблюдения от начала до текущей позиции. Размер окна равен номеру текущего отсчёта плюс единица.

С точки зрения теории вероятностей, это прямое следствие \textbf{закона больших чисел} (strong law of large numbers). Если отсчёты $x[0], x[1], \ldots, x[i]$ --- независимые одинаково распределённые случайные величины с математическим ожиданием $\mu$, то их среднее almost surely сходится к $\mu$ при $i \to \infty$.

\subsection{Скорость сходимости}

\textbf{Центральная предельная теорема} даёт асимптотическое распределение оценки: при большом $i$ величина
\[
\bigl(y[i] - \mu\bigr) \cdot \frac{\sqrt{i}}{\sigma}
\]
стремится к стандартному нормальному распределению $\mathcal{N}(0,1)$. Стандартная ошибка среднего убывает как $\sigma / \sqrt{i}$, что является оптимальной скоростью сходимости для независимых наблюдений.

Дисперсия оценки равна $\sigma^2 / (i+1)$, что монотонно убывает по мере накопления данных. Это делает прогрессивное среднее наиболее эффективным методом для конца сигнала, но крайне неустойчивым для его начала.

\subsection{Несмещённость и состоятельность}

Прогрессивное среднее является \textbf{несмещённой} оценкой постоянного уровня $\mu$: $\mathbb{E}[y[i]] = \mu$ для всех $i$. Более того, оно является \textbf{состоятельной} (consistent) оценкой: при $i \to \infty$ величина $y[i]$ сходится к $\mu$ почти наверное и в квадратичном среднем.

Однако прогрессивное среднее не является \textbf{эффективной} оценкой для сигналов с изменяющимся уровнем, поскольку учитывает все предыдущие наблюдения с одинаковым весом. Если $\mu[i]$ изменяется, старые данные вносят систематическую ошибку, которая не убывает с ростом $i$.

\subsection{Практические последствия}

На практике прогрессивное среднее даёт максимальное сглаживание, но ценой длительной задержки стабилизации. На первых десятках отсчётов результат сильно зависит от конкретных реализаций шума. Метод наиболее подходит для ситуаций, когда предполагается, что сигнал в целом постоянен, а шум значительно превышает изменения уровня.

%% === 4. HYBRID ===
\section{Гибридное сглаживание (Hybrid)}

\subsection{Мотивация с точки зрения адаптивных оценок}

Гибридный метод решает фундаментальную проблему нестационарных оценок: ни один из двух предыдущих методов не является оптимальным на всей области определения сигнала. Boxcar хорош для коротких интервалов (малая дисперсия при фиксированном окне), а progressive --- для длинных (максимальная точность при большом числе наблюдений).

Идея гибридного метода заключается в \textbf{адаптивном переключении} между этими двумя оценками по мере накопления данных.

\subsection{Зоны и интерполяция}

Область определения разбивается на три зоны:
\begin{itemize}
  \item \textbf{Зона I} ($0 \leq i < M$): используется boxcar.
  \item \textbf{Зона II} ($M \leq i < 2M$): линейная интерполяция с коэффициентом $\lambda = (i - M) / M$:
  \[
  y[i] = (1 - \lambda) \cdot y_{\text{boxcar}}[i] + \lambda \cdot y_{\text{prog}}[i].
  \]
  \item \textbf{Зона III} ($i \geq 2M$): используется progressive.
\end{itemize}

С точки зрения теории оценок, линейная интерполяция между двумя несмещёнными оценками также является \textbf{несмещённой оценкой}. Непрерывность перехода гарантирована тем, что на границах зон $\lambda$ принимает значения $0$ и $1$ соответственно.

\subsection{Компромисс смещения и дисперсии}

Гибридный метод минимизирует суммарную среднеквадратичную ошибку (MSE) на всём интервале, сочетая низкую дисперсию boxcar в начале и низкое смещение progressive в конце. В переходной зоне MSE плавно убывает от boxcar-уровня к progressive-уровню.

Такой подход аналогичен \textbf{адаптивным методам оценки} в непараметрической статистике, где размер окна подбирается в зависимости от локальных свойств сигнала. В гибридном методе вместо изменения размера окна происходит изменение типа оценки.

\subsection{Непрерывность и гладкость}

Важным свойством гибридного метода является гарантия \textbf{непрерывности} результата на всех границах зон. В точке $i = M$ значение совпадает с $y_{\text{boxcar}}[M]$, в точке $i = 2M$ --- с $y_{\text{prog}}[2M]$. Линейная интерполяция обеспечивает непрерывность первой производной в переходных точках.

%% === ЗАКЛЮЧЕНИЕ ===
\section{Заключение}

Рассмотренные методы сглаживания образуют последовательность от простейшего (boxcar) к более сложному и адаптивному (hybrid). Каждый из них может быть интерпретирован как особый случай ядерной оценки условного среднего, отличающийся выбором ядра и стратегией определения ширины окна.

\begin{itemize}
  \item \textbf{Boxcar} использует равномерное ядро с фиксированной шириной.
  \item \textbf{Gaussian} заменяет равномерное ядро гауссовым, устраняя артефакты среза.
  \item \textbf{Progressive} увеличивает ширину окна линейно, используя все доступные данные.
  \item \textbf{Hybrid} комбинирует boxcar и progressive, обеспечивая оптимальный компромисс.
\end{itemize}

Все четыре метода являются несмещёнными для постоянного сигнала и состоятельными при стремлении числа наблюдений к бесконечности. Выбор между ними определяется конкретной задачей: степенью шума, характером изменения сигнала и требованиями к гладкости результата.

\end{document}
